\documentclass[a4paper]{article}

\def\npart {药理学精读}
\def\nterm {PK/PD}
\def\nyear {2026}
\def\nlecturer {Holford and Sheiner; Rowland and Tozer}
\def\ncourse {华法林 PK/PD 精读笔记样例}
\def\ncoursehead {华法林 PK/PD}
\def\nauthor {NanFangHong}

\RequirePackage{etex}
\makeatletter
\ifx \npart \undefined
  \def\npart{PK/PD Notes}
\fi
\ifx \nterm \undefined
  \def\nterm{Reading notes}
\fi
\ifx \nyear \undefined
  \def\nyear{\the\year}
\fi
\ifx \nlecturer \undefined
  \def\nlecturer{the source literature}
\fi
\ifx \nauthor \undefined
  \def\nauthor{NanFangHong}
\fi
\ifx \ncourse \undefined
  \def\ncourse{Pharmacology Reading Note}
\fi
\ifx \ncoursehead \undefined
  \def\ncoursehead{\ncourse}
\fi

\author{Based on \nlecturer \\\small Notes taken by \nauthor}
\date{\nterm\ \nyear}
\title{\npart\ --- \ncourse}

\usepackage{amsfonts}
\usepackage{amsmath}
\usepackage{amssymb}
\usepackage{amsthm}
\usepackage{booktabs}
\usepackage{caption}
\usepackage{enumitem}
\usepackage{fancyhdr}
\usepackage{fontspec}
\usepackage{graphicx}
\usepackage{iftex}
\usepackage{mathtools}
\usepackage{microtype}
\usepackage{siunitx}
\usepackage{tabularx}
\usepackage{tikz}
\usepackage{xcolor}
\usepackage{xurl}
\usepackage{xeCJK}

\IfFontExistsTF{Latin Modern Roman}{
  \setmainfont[Ligatures=NoCommon]{Latin Modern Roman}
}{}

\IfFontExistsTF{Noto Serif CJK SC}{
  \setCJKmainfont{Noto Serif CJK SC}
}{
  \IfFontExistsTF{Songti SC}{
    \setCJKmainfont{Songti SC}
  }{
    \IfFontExistsTF{PingFang SC}{
      \setCJKmainfont{PingFang SC}
    }{}
  }
}

\usepackage[hidelinks,unicode,breaklinks=true]{hyperref}
\hypersetup{
  pdfauthor={\nauthor},
  pdfsubject={\npart\ - \ncourse},
  pdftitle={\npart\ - \ncourse},
  pdfkeywords={PK PD pharmacology reading notes \ncourse}
}

\pagestyle{fancyplain}
\lhead{\emph{\nouppercase{\leftmark}}}
\rhead{
  \ifnum\thepage=1
  \else
    \npart\ \ncoursehead
  \fi}

\theoremstyle{definition}
\newtheorem*{aim}{Aim}
\newtheorem*{assumption}{Assumption}
\newtheorem*{defi}{Definition}
\newtheorem*{eg}{Example}
\newtheorem*{notation}{Notation}
\newtheorem*{question}{Question}
\newtheorem*{remark}{Remark}
\newtheorem*{warning}{Warning}

\theoremstyle{plain}
\newtheorem{nthm}{Theorem}[section]
\newtheorem{nprop}[nthm]{Proposition}

\renewcommand{\labelitemi}{--}
\renewcommand{\labelitemii}{$\circ$}
\renewcommand{\labelenumi}{(\roman{*})}

\let\stdsection\section
\renewcommand\section{\newpage\stdsection}

\newcommand{\dd}{\mathop{}\!\mathrm{d}}
\newcommand{\Emax}{E_{\max}}
\newcommand{\pksym}[1]{\ifmmode\text{\textit{#1}}\else\textit{#1}\fi}
\newcommand{\pklab}[1]{\mathrm{#1}}
\newcommand{\ECfifty}{\pksym{EC}_{50}}
\newcommand{\term}[1]{\emph{#1}}
\makeatother


\begin{document}
\maketitle
{\small
  \noindent\textbf{阅读目标}\\
  用一篇小型华法林笔记测试 dec41-style 的完整发布链路：普通 PDF、带目录的 full HTML、按章节切块的 HTML，以及从目录跳回正文位置的链接。\hspace*{\fill}[pipeline]

  \vspace{5pt}
  \noindent\textbf{药理学问题}\\
  这份样例把药动学暴露、药效滞后、效应室模型和浓度--效应曲线放在同一个文档里，方便之后替换成真实文献精读。\hspace*{\fill}[PK/PD]

  \vspace{5pt}
  \noindent\textbf{渲染检查}\\
  目录页用于检查 PDF 内部链接和 pdf2htmlEX full 页面跳转；每个非星号 section 和 subsection 会被切成独立 HTML 页面。\hspace*{\fill}[HTML]
}

\tableofcontents

\section{文献信息}

\subsection{阅读对象}

\begin{tabular}{ll}
\toprule
字段 & 内容 \\
\midrule
主题 & 药动学--药效学关系、效应延迟、浓度--效应建模 \\
样例引用 & Holford and Sheiner, \emph{Kinetics of Pharmacologic Response} \cite{holford1982kinetics} \\
关联教材 & Rowland and Tozer 的 PK/PD 概念框架 \cite{rowland2011clinical} \\
\bottomrule
\end{tabular}

\subsection{精读目的}

这不是对华法林临床用药的完整综述，而是一份工作流样例。真正使用时，我会把文献的研究设计、模型假设、参数估计方法和自己的批注统一放在同一份 \LaTeX{} 源码里。

\begin{remark}
  首页和 PDF 入口服务于浏览；full HTML 服务于快速全文阅读；切块 HTML 服务于像 dec41 那样按章节跳转和引用。
\end{remark}

\section{核心问题}

\subsection{建模问题}

\begin{question}
  作者想解释的是暴露变化、受体/效应系统延迟，还是两者的组合？
\end{question}

\begin{question}
  模型中的状态变量哪些可观测，哪些只能由模型间接推断？
\end{question}

\begin{question}
  参数估计是否足以支持机制解释，还是只支持经验预测？
\end{question}

\subsection{阅读顺序}

\begin{enumerate}[leftmargin=2em]
  \item 先抽取研究问题和给药--采样设计。
  \item 再画出浓度、效应、效应室之间的变量关系。
  \item 最后判断参数是否能由数据识别，还是依赖固定假设。
\end{enumerate}

\section{模型骨架}

\subsection{直接效应模型}

一个最小浓度--效应关系可以写成
\[
  E(C) = \frac{\Emax C}{\ECfifty + C}.
\]
这里 $C$ 是浓度，$E(C)$ 是可观测或可建模的药效，$\Emax$ 是最大效应，$\ECfifty$ 是达到半最大效应所需的浓度。

\begin{assumption}
  如果浓度和效应之间没有明显滞后，直接效应模型可以作为第一层基线模型。
\end{assumption}

\subsection{效应室模型}

如果需要描述药效滞后，可以引入效应室：
\[
  \frac{\dd C_e}{\dd t} = k_{e0}(C_p - C_e),
\]
其中 $C_p$ 是血浆浓度，$C_e$ 是效应室浓度，$k_{e0}$ 控制效应室与血浆之间的平衡速度。

这个模型的阅读重点不是公式本身，而是公式背后的可识别性：如果只有稀疏采样或只有终点效应，$k_{e0}$ 的估计可能会非常依赖先验结构。

\section{精读批注模板}

\subsection{正文批注}

\begin{itemize}[leftmargin=2em]
  \item \textbf{实验设计：}记录给药方式、采样时间、终点指标和受试群体。
  \item \textbf{关键图表：}标出浓度--时间、效应--时间、浓度--效应滞后环。
  \item \textbf{可识别性：}区分由数据支撑的参数和由先验/结构假设固定的参数。
  \item \textbf{我的判断：}说明这篇文献对自己的模型、实验或综述有什么直接帮助。
\end{itemize}

\subsection{复现记录}

每次复现模型时，都应保留输入数据单位、时间单位、初始条件和参数约束。以后同步到个人主页时，这一节可以自然扩展成代码仓库、图表和补充材料的入口。

\section{下一步}

把这份文件复制成新的 \texttt{.tex} 文件，更换标题、引用键和批注内容；同时在 \texttt{notes/notes.json} 里新增一条主页摘要。模板本身由 \texttt{notes/header.tex} 固定，普通 PDF 和 HTML 切块都沿用同一套 dec41-style 入口。

\bibliographystyle{plain}
\bibliography{../refs/pharmacology}

\end{document}
